\chapter{Schermcodes}

\section{Schermcodes}

\label{appendix:screencodes}

Een tekstteken wordt in het schermgeheugen weergegeven door een schermcode. Er bestaan 256 mogelijke schermcodes, die elk verwijzen naar een afbeelding in de huidige tekenset.

Een volledige tekenset bevat twee groepen van 256 afbeeldingen: één voor de hoofdlettermodus en één voor de kleinelettermodus, derhalve 512 afbeeldingen in totaal. Slechts één modus kan tegelijk worden weergegeven. De ingebouwde tekensets gebruiken de eerste 128 tekens van elke groep voor normale tekens en de volgende 128 voor de geïnverteerde versies van dezelfde tekens.

In BASIC biedt de speciale array {\bf T@\&()} toegang tot de tekens op het scherm via kolom- en rij-indexen. De waarden in deze speciale array zijn schermcodes. Het {\bf FONT}-commando schakelt tussen de ingebouwde tekensets. Het {\bf CHARDEF}-commando wijzigt de afbeelding die bij een schermcode hoort.

\underline{OPMERKING}: Schermcodes verschillen van PETSCII-codes. PETSCII-codes worden gebruikt om tekstgegevens op te slaan, te verzenden en te ontvangen en om te bepalen hoe strings op het scherm worden weergegeven. Wanneer een PETSCII-teken op het scherm wordt weergegeven, wordt de bijbehorende schermcode in het schermgeheugen geschreven. Voor een overzicht van PETSCII-codes, zie bijlage \vref{appendix:petsciicodes}.

De volgende tabel geeft een overzicht van de schermcodes. Indien een code afhankelijk van de modus een ander teken oplevert, wordt het teken vermeld als ``hoofdletters / kleine letters''.

\index{Toetsenbord!Schermcodes}
\begin{adjustwidth}{}{-2cm}
\begin{multicols}{4}
\begin{description}[align=left,labelwidth=0.2cm]
    \item [0]   @    \item [1]   A / a    \item [2]   B / b    \item [3]   C / c    \item [4]   D / d    \item [5]   E / e    \item [6]   F / f    \item [7]   G / g    \item [8]   H / h    \item [9]   I / i    \item [10]  J / j    \item [11]  K / k    \item [12]  L / l    \item [13]  M / m    \item [14]  N / n    \item [15]  O / o    \item [16]  P / p    \item [17]  Q / q    \item [18]  R / r    \item [19]  S / s    \item [20]  T / t    \item [21]  U / u    \item [22]  V / v    \item [23]  W / w    \item [24]  X / x    \item [25]  Y / y    \item [26]  Z / z    \item [27]  [    \item [28]  \pounds    \item [29]  ]    \item [30]  $\uparrow$    \item [31]  $\leftarrow$    \item [32]  spatie    \item [33]  !    \item [34]  "    \item [35]  \#    \item [36]  \$    \item [37]  \%    \item [38]  \&    \item [39]  '    \item [40]  (    \item [41]  )    \item [42]  *    \item [43]  +    \item [44]  ,    \item [45]  -    \item [46]  .    \item [47]  /    \item [48]  0    \item [49]  1    \item [50]  2    \item [51]  3    \item [52]  4    \item [53]  5    \item [54]  6    \item [55]  7    \item [56]  8    \item [57]  9    \item [58]  :    \item [59]  ;    \item [60]  <    \item [61]  =    \item [62]  >    \item [63]  ?    \item [64]  \graphicsymbol{C}    \item [65]  \graphicsymbol{A} / A    \item [66]  \graphicsymbol{B} / B    \item [67]  \graphicsymbol{C} / C    \item [68]  \graphicsymbol{D} / D    \item [69]  \graphicsymbol{E} / E    \item [70]  \graphicsymbol{F} / F    \item [71]  \graphicsymbol{G} / G    \item [72]  \graphicsymbol{H} / H    \item [73]  \graphicsymbol{I} / I    \item [74]  \graphicsymbol{J} / J    \item [75]  \graphicsymbol{K} / K    \item [76]  \graphicsymbol{L} / L    \item [77]  \graphicsymbol{M} / M    \item [78]  \graphicsymbol{N} / N    \item [79]  \graphicsymbol{O} / O    \item [80]  \graphicsymbol{P} / P    \item [81]  \graphicsymbol{Q} / Q    \item [82]  \graphicsymbol{R} / R    \item [83]  \graphicsymbol{S} / S    \item [84]  \graphicsymbol{T} / T    \item [85]  \graphicsymbol{U} / U    \item [86]  \graphicsymbol{V} / V    \item [87]  \graphicsymbol{W} / W    \item [88]  \graphicsymbol{X} / X    \item [89]  \graphicsymbol{Y} / Y    \item [90]  \graphicsymbol{Z} / Z    \item [91]  \graphicsymbol{+}    \item [92]  \graphicsymbol{-}    \item [93]  \graphicsymbol{B}    \item [94]  \graphicsymbol{3} / \graphicsymbol{>} % pi and inverted checker    \item [95]  \graphicsymbol{]} / \graphicsymbol{:} % upper right triangle and back slashes    \item [96]  spatie    \item [97]  \graphicsymbol{j}    \item [98]  \graphicsymbol{i}    \item [99]  \graphicsymbol{t}    \item [100] \graphicsymbol{[}    \item [101] \graphicsymbol{g}    \item [102] \graphicsymbol{<}    \item [103] \graphicsymbol{m}    \item [104] \graphicsymbol{/}    \item [105] \graphicsymbol{?} / \graphicsymbol{)} % upper-left triangle and forward slashes    \item [106] \graphicsymbol{n}    \item [107] \graphicsymbol{q}    \item [108] \graphicsymbol{d}    \item [109] \graphicsymbol{z}    \item [110] \graphicsymbol{s}    \item [111] \graphicsymbol{p}    \item [112] \graphicsymbol{a}    \item [113] \graphicsymbol{e}    \item [114] \graphicsymbol{r}    \item [115] \graphicsymbol{w}    \item [116] \graphicsymbol{h}    \item [117] \graphicsymbol{j}    \item [118] \graphicsymbol{l}    \item [119] \graphicsymbol{y}    \item [120] \graphicsymbol{u}    \item [121] \graphicsymbol{p}    \item [122] \graphicsymbol{\{} / \graphicsymbol{8} % lower right square edge and tick symbol    \item [123] \graphicsymbol{f}    \item [124] \graphicsymbol{c}    \item [125] \graphicsymbol{x}    \item [126] \graphicsymbol{v}    \item [127] \graphicsymbol{b}
\end{description}
\end{multicols}
\end{adjustwidth}

\underline{OPMERKING}: In de ingebouwde tekensets zijn codes 128-255 geïnverteerde versies van codes 0-127.
